%% kapa: exam=1 round=37 year=2026 period=2 form=A subject=accounting subject_ko=회계학 seq=31 number=71
%% source: Q-Net 감정평가사 기출문제 (cst00309) artlSeq=5258924 file=raw/1008/5258924/2026년 제37회 감정평가사 1차 2교시.pdf
%% conversion: kordoc 4.24.51 -> tools/gen1.py (auto); review=proofread-sonnet; flags=none
\begin{question}{71}
최신의 관리회계 기법에 관한 설명으로 옳지 않은 것은?
\begin{choices}
  \item 품질원가에서 설계품질이란 설계 사양에 맞게 제품을 생산하였는지로 측정하는 것을 말한다.
  \item 수명주기원가계산에서는 제조이전 단계를 철저히 통제함으로써 원가 발생을 줄일 수 있다.
  \item 카이젠원가는 제조단계에서 원가를 절감하는 방법으로 일반적으로 원가절감의 여지가 크지 않다.
  \item 원가기획은 제조원가의 대부분이 결정되는 제품의 기획ㆍ설계 단계에서부터 원가절감을 위한 조치가 가능하다.
  \item 균형성과표는 재무적 관점, 고객 관점, 내부프로세스 관점, 학습과 성장 관점의 네 가지 관점을 반영한 균형 잡힌 성과관리제도를 말한다.
\end{choices}
\end{question}
\answer{1}
